\chapter{Reproduction and Validation Material}
\label{chap:appendices}

\section{Fresh-Start Command-Prompt Procedure}
\label{app:fresh-start}

The following procedure uses Windows Command Prompt. It does not require
PowerShell.

\begin{enumerate}
    \item Open Command Prompt.
    \item Change to the repository:
\begin{lstlisting}[language={},caption={Open the NEMO repository}]
cd /d C:\Users\muigu\Documents\Projects\NEMO
\end{lstlisting}
    \item Create or update the project environment:
\begin{lstlisting}[language={},caption={Set up the project}]
build.bat setup
\end{lstlisting}
    \item Run the offline checks:
\begin{lstlisting}[language={},caption={Run offline verification}]
build.bat check
\end{lstlisting}
    \item List the registered parts:
\begin{lstlisting}[language={},caption={List the registered components}]
.venv\Scripts\nemo.exe parts
\end{lstlisting}
    \item Evaluate the three baselines:
\begin{lstlisting}[language={},caption={Evaluate the analytical baselines}]
.venv\Scripts\nemo.exe evaluate --part bracket
.venv\Scripts\nemo.exe evaluate --part padeye
.venv\Scripts\nemo.exe evaluate --part stabilizer
\end{lstlisting}
\end{enumerate}

\section{Bracket Proof-of-Concept Procedure}
\label{app:bracket-poc}

\begin{enumerate}
    \item Start Autodesk Fusion.
    \item Open \emph{Utilities}, then \emph{Scripts and Add-Ins}.
    \item Select and start NEMOBridge.
    \item If a bracket opens automatically, leave it available. NEMOBridge can
    load or create the demonstration design during startup.
    \item Open Command Prompt and change to the repository.
    \item Run:
\begin{lstlisting}[language={},caption={Run the preserved bracket proof of concept}]
cd /d C:\Users\muigu\Documents\Projects\NEMO
.venv\Scripts\python.exe demos\bracket-proof-of-concept\manual_fusion_handshake.py
\end{lstlisting}
    \item Return to Fusion and inspect the regenerated bracket.
    \item For the complete manual study, follow
    Section~\ref{sec:bracket-poc-method}.
\end{enumerate}

\section{Schema-Version-2 Request Example}
\label{app:schema-request}

\begin{lstlisting}[language=Python,caption={Illustrative bracket CAD request},
label={lst:request-example}]
{
  "schema_version": 2,
  "part_id": "bracket",
  "run_id": "demo_bracket",
  "iteration": 0,
  "mode": "fusion_cad",
  "parameters": {
    "baseplate_length": 150.0,
    "baseplate_width": 100.0,
    "baseplate_thickness": 8.0,
    "rib_height": 40.0,
    "rib_thickness": 6.0,
    "fillet_radius": 5.0
  },
  "artifact_formats": ["step", "boundary_tags"]
}
\end{lstlisting}

\section{CAD-Only Response Example}
\label{app:schema-response}

\begin{lstlisting}[language=Python,caption={Illustrative partial CAD response},
label={lst:response-example}]
{
  "schema_version": 2,
  "part_id": "bracket",
  "run_id": "demo_bracket",
  "iteration": 0,
  "status": "partial",
  "metrics": {
    "mass_kg": 0.4228127502,
    "volume_m3": 0.0001565973,
    "max_stress_mpa": null,
    "factor_of_safety": null,
    "max_deflection_mm": null,
    "objective_value": null
  },
  "artifacts": {
    "step": "...\\bracket_0000.step",
    "boundary_tags": "...\\bracket_0000.boundaries.json"
  },
  "error": "Native CAD generation and mass extraction completed.
            Static-stress results require manual Fusion validation."
}
\end{lstlisting}

The partial status is intentional. It prevents a CAD success from being
interpreted as an FEA success.

\section{Selected Finalist Parameters}
\label{app:finalist-parameters}

\begin{table}[H]
\centering
\caption{Selected bracket finalist parameters}
\label{tab:app-bracket-finalist}
\begin{tabularx}{\textwidth}{Xr}
\toprule
\textbf{Parameter} & \textbf{Value (mm)}\\
\midrule
Baseplate length & 100.0328\\
Baseplate width & 80.9380\\
Baseplate thickness & 4.0429\\
Rib height & 33.2576\\
Rib thickness & 3.0000\\
Fillet radius & 3.5308\\
\bottomrule
\end{tabularx}
\end{table}

\begin{table}[H]
\centering
\caption{Selected padeye finalist parameters}
\label{tab:app-padeye-finalist}
\begin{tabularx}{\textwidth}{Xr}
\toprule
\textbf{Parameter} & \textbf{Value (mm)}\\
\midrule
Base length & 220.3878\\
Base width & 150.4654\\
Base thickness & 8.0484\\
Lug height & 141.7256\\
Lug thickness & 12.1613\\
Neck width & 80.0000\\
Gusset height & 70.0000\\
Gusset thickness & 6.2101\\
Fillet radius & 25.0000\\
\bottomrule
\end{tabularx}
\end{table}

\begin{table}[H]
\centering
\caption{Selected stabilizer finalist parameters}
\label{tab:app-stabilizer-finalist}
\begin{tabularx}{\textwidth}{Xlr}
\toprule
\textbf{Parameter} & \textbf{Unit} & \textbf{Value}\\
\midrule
Skin thickness & mm & 4.0106\\
Front spar position & \% chord & 25.0000\\
Rear spar position & \% chord & 72.0000\\
Front spar thickness & mm & 7.8522\\
Rear spar thickness & mm & 6.5297\\
Rib thickness & mm & 5.6971\\
Root insert length & mm & 221.4762\\
Root insert thickness & mm & 13.6624\\
Root fillet radius & mm & 18.3273\\
\bottomrule
\end{tabularx}
\end{table}

\section{Manual FEA Validation Record}
\label{app:fea-record}

Each baseline and finalist FEA record should contain the fields in
Table~\ref{tab:app-fea-checklist}.

\begin{longtable}{p{0.25\textwidth}p{0.65\textwidth}}
\caption{Minimum manual FEA validation checklist}
\label{tab:app-fea-checklist}\\
\toprule
\textbf{Record item} & \textbf{Required evidence}\\
\midrule
\endfirsthead
\toprule
\textbf{Record item} & \textbf{Required evidence}\\
\midrule
\endhead
\bottomrule
\endfoot
Candidate identity &
Part identifier, candidate identifier, source run, iteration, parameter table,
schema version, and generator version.\\

Geometry &
Native CAD screenshot, connected-body count, volume, mass, and STEP path.\\

Material &
Name, density, Young's modulus, Poisson ratio, yield strength, and source.\\

Load &
Magnitude, unit, derivation, direction, selected faces, distribution, and
force-per-entity state.\\

Support &
Constraint type, selected faces, physical interpretation, and screenshot.\\

Contact &
Contact type, interfaces, assumptions, and any nonlinear settings.\\

Mesh &
Element type and order, global size, local sizes, growth control, node count,
element count, and screenshot.\\

Equilibrium &
Applied-load vector, support-reaction vector, imbalance, and percentage
difference.\\

Stress &
Maximum von Mises stress, location, stable nearby monitored stress, and
contour screenshot.\\

Factor of safety &
Minimum value, location, yield definition, and contour screenshot.\\

Displacement &
Maximum magnitude, monitored displacement, direction, and location.\\

Convergence &
At least three meshes, same monitored locations, result changes, and selected
mesh justification.\\

Acceptance &
FOS and displacement pass/fail, local failure-mode checks, limitations, and
engineering disposition.\\
\end{longtable}

\section{Bracket Corrected-Study Record Template}
\label{app:bracket-study-template}

\begin{table}[H]
\centering
\caption{Bracket Static Stress input record}
\label{tab:app-bracket-input}
\begin{tabularx}{\textwidth}{p{0.30\textwidth}X}
\toprule
\textbf{Field} & \textbf{Required entry}\\
\midrule
Material & Project Aluminum 6061-T6 properties\\
Analytical load & 1471.5~N\\
Fusion force input & \textbf{1472 N total}\\
Force per entity & Off\\
Loaded faces & Four upper sloping rib faces\\
Load direction & Downward in the declared study coordinate system\\
Supports & Four mounting-hole cylindrical faces\\
Element order & Parabolic\\
Local refinements & Mounting holes and rib-root fillets\\
Reaction target & Approximately 1472~N opposite the applied direction\\
Constraint checks & FOS $\geq2.5$ and displacement $\leq0.5$~mm\\
\bottomrule
\end{tabularx}
\end{table}

\section{Mesh-Convergence Template}
\label{app:mesh-template}

\begin{table}[H]
\centering
\caption{Mesh-convergence recording template}
\label{tab:app-mesh-template}
\small
\resizebox{\textwidth}{!}{%
\begin{tabular}{crrrrr}
\toprule
\textbf{Mesh} & \textbf{Nodes} & \textbf{Elements} &
\textbf{Peak stress (MPa)} & \textbf{Monitored stress (MPa)} &
\textbf{Displacement (mm)}\\
\midrule
1 & & & & &\\
2 & & & & &\\
3 & & & & &\\
4 & & & & &\\
\bottomrule
\end{tabular}%
}
\end{table}

Mesh 1 is the coarse reference, meshes 2 and 3 are successive refinements, and
mesh 4 is an optional confirmation. The monitored stress must use the same
geometric region on every mesh. A moving global maximum is not a valid
convergence quantity when it moves between unrelated singular edges.

\section{Stabilizer Pressure-Band Calculation}
\label{app:pressure-bands}

Let the unnormalized weights be

\begin{equation}
\mathbf{w}=[1.00,\ 0.93,\ 0.75,\ 0.40].
\end{equation}

Their sum is 3.08. The band-force fractions are:

\begin{equation}
\hat{\mathbf{w}}=
[0.324675,\ 0.301948,\ 0.243506,\ 0.129870].
\end{equation}

For the 11\,808~N resultant, the approximate band forces are:

\begin{table}[H]
\centering
\caption{Prescribed stabilizer pressure-band forces}
\label{tab:app-pressure-bands}
\begin{tabular}{crr}
\toprule
\textbf{Band} & \textbf{Normalized fraction} & \textbf{Force (N)}\\
\midrule
1 & 0.324675 & 3833.8\\
2 & 0.301948 & 3565.4\\
3 & 0.243506 & 2875.1\\
4 & 0.129870 & 1533.8\\
\midrule
Total & 1.000000 & 11\,808.0\\
\bottomrule
\end{tabular}
\end{table}

The FEA pressure for a band is its force divided by its selected CAD face area.
The forces are prescribed engineering inputs and are not CFD results.

\section{Evidence Inventory}
\label{app:evidence-inventory}

\begin{longtable}{>{\raggedright\arraybackslash}p{0.34\textwidth}p{0.56\textwidth}}
\caption{Principal repository evidence used in the report}
\label{tab:app-evidence}\\
\toprule
\textbf{Location} & \textbf{Evidence}\\
\midrule
\endfirsthead
\toprule
\textbf{Location} & \textbf{Evidence}\\
\midrule
\endhead
\bottomrule
\endfoot
\path{src/nemo/parts/definitions} &
Authoritative part parameters, units, materials, loads, constraints, fixed
geometry, and boundary roles.\\

\path{src/nemo/parts/analytical.py} &
First-order volume, stress, factor-of-safety, and displacement equations.\\

\path{src/nemo/optimizer.py} &
Scaled bounded Nelder--Mead search and evaluation handling.\\

\path{src/nemo/logger.py} &
Stable one-part-per-log CSV and JSON records.\\

\path{fusion_addin/NEMOBridge} &
Fusion add-in, native generators, artifact export, and boundary selectors.\\

\path{reports/final_20260728_*_validation} &
Baseline and finalist parameters, analytical metrics, Fusion requests, Fusion
responses, and checklists.\\

\path{F360 report/Study Final.html} &
Bracket manual study setup, mesh metadata, result views, and the recorded
1.472~N force error.\\

\path{demos/bracket-proof-of-concept} &
Preserved initial demonstration flow and historical validation material.\\

\path{reports/project-report} &
LaTeX report source, bibliography, extracted images, reproducible chart
script, and compiled PDF.\\
\end{longtable}

\section{Demonstration Sequence}
\label{app:demo-sequence}

A concise defense demonstration can follow this order:

\begin{enumerate}
    \item state the problem and evidence boundary;
    \item show the three part definitions;
    \item evaluate the three analytical baselines;
    \item show the 60-sample and three-start workflow;
    \item compare baseline and finalist mass;
    \item open one stored validation package;
    \item send or replay a bracket CAD request;
    \item show the native bracket, mass, STEP, and boundary metadata;
    \item identify the four support faces and four upper load faces;
    \item explain the 1.472~N error and the corrected 1472~N procedure;
    \item show how padeye and stabilizer roles extend the same architecture;
    \item close with the validation work required before manufacture.
\end{enumerate}

\section{Report Metadata Requiring Confirmation}
\label{app:metadata-placeholders}

The title-page project code remains \code{MR-PRO-26-XX}, and the supervisor is
shown as \code{SUPERVISOR NAME}. These placeholders must be replaced with the
official project code and supervisor name before final institutional
submission.
